\vspace{-0.2cm}
\section{Introduction}
\vspace{-0.1cm}
Using LLMs as automated judges has become standard practice for evaluating model outputs~\citep{zheng2023judging,tan2025judgebench,li2024arenahard}. On text-centric tasks -- dialogue, summarization, instruction following -- judge reliability is well characterized, with documented biases and known failure modes~\citep{zheng2023judging,wang2023fair}. As LLMs are increasingly deployed as autonomous agents that invoke tools and orchestrate multi-step workflows, the judge paradigm has been extended naturally to agentic tool-calling~\citep{qin2024toolllm,guo2024stabletoolbench,guo2025mcpagentbench}. This extension, however, has proceeded without a basic calibration check: \emph{how reliable are LLM judges in this structured setting?}

Agentic tool-calling differs from text evaluation in ways that matter for judge reliability. Correctness is not a matter of fluency or preference: it requires selecting the right tools from a typed schema, supplying well-formed arguments, ordering calls to respect execution dependencies, and covering all parts of the user's intent. A plan can fail in four orthogonal ways (tool selection, parameter structure, sequence accuracy, query coverage) that do not correlate cleanly -- a judge reliable at one may be blind to another, e.g., calibrated on simple sequential tasks yet fail on complex fan-in or diamond workflows where parallel branches must converge. When no ground-truth execution trace is available, the common deployment scenario, the judge must instead reconstruct correctness from the query and tool schemas alone, a fundamentally harder inference problem. Existing benchmarks that deploy LLM judges for tool-calling report only aggregate pass-rate agreement and vary none of these dimensions~\citep{qin2024toolllm,guo2024stabletoolbench,guo2025mcpagentbench}, leaving practitioners without a principled basis for choosing a judge, setting its configuration, or interpreting its outputs.

We introduce \textbf{AgentJudgeBench} to close this gap. The benchmark comprises 3,808 BFCL-style records~\citep{patil2025bfcl} spanning six DAG\footnote{A Directed Acyclic Graph (DAG) encodes execution dependencies between tool calls: nodes are tool invocations and directed edges indicate that one call must complete before another can begin, with the acyclicity constraint preventing circular dependencies.} topologies at three controlled difficulty tiers, each with a programmatically verified ground-truth trace. Five generators -- four open-weight (3B--70B) and one frontier (GPT-5.4) -- produce tool-calling outputs that are scored by six LLM judges (ranging from 20B open-weight to frontier closed systems) under paired with- and without-GT conditions across four structural metrics (tool selection, parameter structure, sequence accuracy, query coverage), yielding $321{,}648$ paired with-GT/without-GT evaluations (Appendix~\ref{app:eval-stats}). We organize our analysis around six research questions: which metrics and topologies are hardest (RQ1), how much judges agree with each other (RQ2), why without-GT alignment converges on hard queries (RQ3), and how sensitive alignment is to judge temperature (RQ4), chain-of-thought reasoning (RQ5), and prompt format (RQ6).

The results are non-obvious. Judge alignment degrades monotonically with query difficulty, $1.5\times$ faster without ground truth than with it, and all six judges converge to a $77$--$82\%$ ceiling on hard queries without a reference, regardless of model capacity; a follow-up ablation (\S\ref{sec:exp:ceiling}) confirms task difficulty as the primary driver for capable generators, though its height is somewhat prompt-dependent for weaker ones. GT exposure is not universally beneficial: two frontier judges (GPT-5.4, Gemini-2.5-Pro) are \emph{less} aligned when shown the reference, consistent with over-anchoring rather than independent judgement. Chain-of-thought adds at most $0.3$\,pp across 24 paired comparisons, and temperature has negligible effect (${\leq}0.25$\,pp spread, two judge-generator pairings). Structured per-metric prompt rubrics add $+4.8$--$+6.5$\,pp over free-form on one pairing, the largest lever we test, but a second pairing shows a smaller effect that reverses on hard queries, so we treat prompt format as impactful but judge/generator-dependent rather than universally dominant (\S\ref{sec:abl:format}). Inter-judge agreement is moderate ($\kappa \approx 0.42$ with GT), with systematic, capacity-correlated patterns not visible in aggregate scores.

\noindent\textbf{Scope.} We study LLM-as-a-judge specifically for agentic tool-calling, not the broader LLM-as-a-judge literature; DAG-structured tool-use data, synthetic dependency graphs, and programmatic trajectory scoring each have close prior work (\S2). Our contribution is the reliability \emph{protocol} that combines them -- paired with-GT/without-GT conditions, controlled difficulty, and per-metric decomposition -- applied to measure judge reliability rather than agent capability.

\noindent\textbf{Contributions.} \textbf{(1)}~A dataset of 3,808 records spanning six DAG topologies and three difficulty tiers with programmatically verified ground-truth traces, to be released publicly. \textbf{(2)}~A four-metric evaluation framework with a paired with-GT/without-GT protocol and bootstrap confidence intervals. \textbf{(3)}~A systematic six-RQ empirical study yielding actionable guidance: with ground truth, QwQ-32B best matches the programmatic reference, while GPT-OSS-120B is the most human-aligned judge in our validation study (Appendix~\ref{app:prog-judge-validation}); without ground truth, frontier judges lead narrowly, but the convergence ceiling limits the practical difference.
